\chapter{Prepare the Store Rails Toolchain}

The blank store jail can reach the catalog API. We can now add Ruby and Rails
without confusing an application dependency failure with a jail-networking
failure. The store uses the same small package boundary already proven in the
catalog jail.

\section{Exercise 10A: Install Ruby and Rails}

Run the following commands as \texttt{root} in the store jail:

\begin{lstlisting}[language=sh]
pkg install -y ruby rubygem-bundler rubygem-pg \
  rubygem-nokogiri
gem install rails --no-document
\end{lstlisting}

FreeBSD supplies Ruby and Bundler. It also supplies the native PostgreSQL and
Nokogiri extensions. Rails comes from RubyGems. The \texttt{--no-document}
option avoids generating unneeded RI and RDoc files inside the jail. Exact
patch versions are deliberately not part of this exercise; the goal is a
working application stack, not matching the catalog jail gem for gem.

Verify the two capabilities the store application needs:

\begin{lstlisting}[language=sh]
rails --version
ruby -rpg -e 'puts PG.library_version'
\end{lstlisting}

The first command should print the installed Rails release. The second should
print a PostgreSQL client-library version without raising an exception.

\section{Checkpoint}

Stop when both verification commands succeed. Do not generate the application
yet. The next exercise will create a dedicated \texttt{store} account and a
browser-facing Rails application owned by that account.
